\section{Proofs for Section~\ref{sec:fresh}: Fresh Tuples}
\label{app:fresh}

Throughout, $W=(w_1,\dots,w_k)$ is a fresh tuple of a table with $k$ starts
and $N$ states, and $q_l(c)$ is the state of coordinate $l$ at the cut $c$.

\subsection{Width}\label{app:fresh-width}

\begin{lemma}[Width]\label{lem:width}
\begin{enumerate}
\item\label{it:width-word} For every text $z$ and every coordinate $l$, at
  most $N-1$ cuts $b<\len{w_l}$ satisfy $\wsuf{w_l}b\pre z$.
\item\label{it:width-all} For every text $z$, the set $\operatorname{Pre}(z)$
  of the numbers $p$ with $1\le p\le\len z$ such that $\wpre zp$ is a suffix
  of some word $w_l$ has at most $k(N-1)$ elements.
\end{enumerate}
\end{lemma}

\begin{proof}
\eqref{it:width-word} The rests $\wsuf{w_l}b$ in question are prefixes of
one text, hence pairwise comparable. By \cref{lem:distinct}, the states at
their cuts are pairwise different, and they differ from the endpoint because
$w_l$ is fresh and $b<\len{w_l}$. \eqref{it:width-all} If $\wpre zp$ is a
suffix of $w_l$, then $b=\len{w_l}-p$ is a cut as in \eqref{it:width-word},
and different $p$ give different $b$ for the same $l$.
\end{proof}

\subsection{Markers}\label{app:fresh-markers}

An anchor of a text $u$ is a cut at which a marker ends that fits into $u$;
an occurrence of a marker is always a complete occurrence.

\begin{proof}[Proof of \cref{lem:markers}]
By induction on $t$ we prove \eqref{it:mk-anti} for level $t$, and
\eqref{it:mk-two}, \eqref{it:mk-nest} for the passage from $t$ to $t+1$. At
level $1$ the markers are single letters, and a letter is a factor of another
letter only if they are equal. If an occurrence of a marker lies inside an
occurrence of another marker of the same level, the first is a factor of the
second, so the texts are equal and the occurrences coincide; this is the
second statement of \eqref{it:mk-anti}. In particular, two markers of one
level that end at the same cut of a text have the same text, so ``the
longest marker ending at $e$'' in \cref{def:markers} is the only one, up to
the choice of the coordinate.

Assume \eqref{it:mk-anti} for level $t$. Let $l\in\Act_{t+1}$, $\ell=\len{w_l}$,
let $e<\ell$ be the last level-$t$ anchor of $w_l$ before its end and $\mu$ the
level-$t$ marker ending at $e$, so that $\mk{t+1}l=\wsuf{w_l}{e-\len\mu}$.
The occurrence of $\mu$ at $[e-\len\mu,e)$ does not lie inside the occurrence
of $\mk tl$ at $[\ell-\len{\mk tl},\ell)$, because the two occurrences end at
different cuts. Hence $e-\len\mu<\ell-\len{\mk tl}$, and $\mk tl$ is a proper
suffix of $\mk{t+1}l$. Now consider a level-$t$ marker $\mu'$ that ends at a
cut $e'$ of the text $\mk{t+1}l$ and fits into it. If $e'\le\len\mu$, its
occurrence lies inside that of the prefix $\mu$, so $e'=\len\mu$. If
$\len\mu<e'<\len{\mk{t+1}l}$, then $w_l$ has a level-$t$ anchor strictly
between $e$ and $\ell$, contradicting the choice of $e$. So the anchors are
$\len\mu$ and the end, which proves \eqref{it:mk-two}. For
\eqref{it:mk-nest}, a level-$(t+1)$ marker that ends at a cut of a text has
the level-$t$ marker of the same coordinate as a suffix, which ends at the
same cut.

For \eqref{it:mk-anti} at level $t+1$, let $\mk{t+1}l$ occur in $\mk{t+1}m$
at the position $o$, and let $\mu_l,\mu_m$ be the level-$t$ markers with
which they begin. The two level-$t$ anchors of the text $\mk{t+1}l$ give
level-$t$ anchors $o+\len{\mu_l}$ and $o+\len{\mk{t+1}l}$ of the text
$\mk{t+1}m$, which by \eqref{it:mk-two} are $\len{\mu_m}$ and
$\len{\mk{t+1}m}$. So the occurrence ends at the end of $\mk{t+1}m$, and
$\mu_l$ and $\mu_m$ end at the same cut of $\mk{t+1}m$; they are equal by
\eqref{it:mk-anti} at level $t$, hence $o=0$ and the two markers are equal.
\end{proof}

We also note that for $t\ge1$ a coordinate $l\notin\Act_{t+1}$ has no
level-$t$ anchor other than its end: if $l\in\Act_t$ this is the definition,
and otherwise $l$ left the hierarchy at a level $r<t$, has no level-$r$
anchor before its end, and level-$t$ anchors are level-$r$ anchors by
\cref{lem:markers}\eqref{it:mk-nest}.

\begin{proof}[Proof of \cref{lem:depth-span}]
\eqref{it:depth} We show by induction that for $l\in\Act_t$ the set
$\operatorname{Pre}(\mk tl)$ of \cref{lem:width} has at least $t$ elements;
then $t\le k(N-1)$ whenever $\Act_t\ne\varnothing$. For $t=1$ the marker is
a suffix of $w_l$. For $t+1$, the marker $\mk{t+1}l$ begins with a level-$t$
marker $\mu$ of some coordinate, whose $t$ prefixes that are suffixes of
words are prefixes of $\mk{t+1}l$ of length at most $\len\mu$, and
$\mk{t+1}l$ itself is a longer one.

\eqref{it:span} For $t=0$ the marker is a letter. Let $t\ge1$, let $\sigma$
be the level-$(t+1)$ marker with the occurrence $\wsub uab$, and let $e$ be
a level-$t$ anchor of $u$ with $a<e\le b$, with a level-$t$ marker $\mu$
ending at $e$. If $\mu$ begins at or after $a$, then $e-a$ is one of the two
anchors of the text $\sigma$. Otherwise $\wsub uae$ is a proper suffix of
$\mu$, hence a suffix of a word, and a prefix of $\sigma$; so
$e-a\in\operatorname{Pre}(\sigma)$, a set of at most $k(N-1)$ elements.
\end{proof}

\subsection{Parses}\label{app:fresh-parses}

Fix a level $t$ between $0$ and the top level. Let $z$ be a text that
occurs in a text $uz$ after the left context $u$. For $t\ge1$, a position
$p$ with $1\le p\le\len z$ is an \emph{inner anchor} of $z$ if a level-$t$
marker of length at most $p$ ends at the cut $p$ of $z$, and a
\emph{crossing anchor} in the context $u$ if a level-$t$ marker of length
greater than $p$ ends at the cut $\len u+p$ of $uz$.

\begin{lemma}[Parsing is local]\label{lem:parse-local}
Let $t\ge1$.
\begin{enumerate}
\item\label{it:pl-union} The level-$t$ anchors of $uz$ after $\len u$ are
  the cuts $\len u+p$ with $p$ an inner anchor of $z$ or a crossing anchor
  in the context $u$. The inner anchors do not depend on $u$.
\item\label{it:pl-order} Every crossing anchor is smaller than every inner
  anchor.
\item\label{it:pl-pre} Every crossing anchor lies in $\operatorname{Pre}(z)$.
\end{enumerate}
Hence the lower anchors of two occurrences of $z$, in arbitrary contexts,
agree from the first inner anchor of $z$ on, and each occurrence has at most
$k(N-1)$ lower anchors before it.
\end{lemma}

\begin{proof}
\eqref{it:pl-union} is the definition. \eqref{it:pl-order} Let a marker
$\mu$ of length greater than $p$ end at $\len u+p$ and a marker $\mu'$ of
length at most $p'$ end at $\len u+p'$, with $p'\le p$. Then the occurrence
of $\mu'$ lies inside that of $\mu$, so by
\cref{lem:markers}\eqref{it:mk-anti} they coincide, which is impossible
since $\mu'$ begins inside $z$ and $\mu$ before it.
\eqref{it:pl-pre} The prefix $\wpre zp$ is a proper suffix of the marker,
which is a suffix of a word.
\end{proof}

For $t=0$ every cut is a lower anchor, in every context.

\begin{lemma}[Texts determine traversal]\label{lem:traverse}
Let $(Z,a,b)$ be an edge at level $t$, let $i$ be a depth with
$a+k(N-1)+2\le i\le b$, and let $z$ be the text of the last $i$ lower blocks
of $Z$. If $z$ ends at an upper anchor $e$ of a word $w_l$, then the upper
occurrence of $w_l$ that ends at $e$ traverses the edge.
\end{lemma}

\begin{proof}
Let $o_1$ be an upper occurrence with parse $Z$; its last $i$ lower blocks
form an occurrence of $z$ whose $i$ lower anchors are the cuts at the depths
$i-1,\dots,0$. Let $o_2$ be the upper occurrence of $w_l$ ending at $e$, and
consider the occurrence of $z$ that ends at $e$.

We claim that $o_2$ begins at or before the first inner anchor $p_0$ of
$z$, measured in this occurrence; for $t=0$ we let $p_0=1$. Suppose an upper
anchor of $w_l$ lies at a position $r$ of $z$ with $0<r<\len z$, and let
$\sigma$ be a level-$(t+1)$ marker ending there. If $\sigma$ fits into $z$,
it also ends at the corresponding cut of $o_1$, which is then an upper
anchor strictly inside the upper block $o_1$, a contradiction; for $t=0$,
where $\sigma$ is a letter, this is the only case. So $\sigma$ begins before
$z$. By \cref{lem:markers}\eqref{it:mk-two}, the text $\sigma$ has no
level-$t$ anchor other than the end of its first level-$t$ marker and its
own end; an inner anchor of $z$ at a position $r'<r$ would be a third one,
since its marker begins after the beginning of $\sigma$ and ends before its
end. So $r\le p_0$, which proves the claim.

For $t=0$ the two parses consist of the letters of $z$. For $t\ge1$, by
\cref{lem:parse-local}, in the occurrence inside $o_1$ at most $k(N-1)$
of the $i$ lower anchors are crossing anchors, so $z$ has at least
$i-k(N-1)$ inner anchors, and in both occurrences the lower anchors from
$p_0$ on are exactly the inner anchors. By the claim, the lower blocks
between them belong to $o_2$. So the parses of $o_1$ and $o_2$ end with the
same $i-k(N-1)-1>a$ lower blocks. Since the parse of $o_1$ is $Z$, the
number $\lcsuf(Y,Z)$ for the parse $Y$ of $o_2$ is a node on the path of $Z$
greater than $a$, and no node lies strictly between $a$ and $b$; so it is at
least $b$.
\end{proof}

\subsection{Block Deletions}\label{app:fresh-deletion}

Fix a block deletion at level $t$, with the edge $(Z,a,b)$ and the depths
$i<j$. Let $z_\lambda$ be the text of the last $i$ lower blocks of $Z$,
$\lambda=\len{z_\lambda}$, and let $F$ be the text of its lower blocks at
the depths $i,\dots,j-1$, $\delta=\len F\ge1$. For a traversing upper
occurrence $o$ of a word $w_l$, between the cuts $s_o<e_o$, let $c_o$ and
$c_o'=c_o+\delta=e_o-\lambda$ be its cuts at the depths $j$ and $i$. The
\emph{site} of $o$ is the interval of positions $[c_o,c_o')$; it carries the
text $F$. The edited tuple $U=(u_1,\dots,u_k)$ is obtained by deleting all
sites.

\subparagraph*{Deleting positions.}
We use the following notation for the deletion of a set of pairwise disjoint
intervals of positions from a word $w$, with the result $u$. A cut of $w$ is
\emph{retained} if it does not lie strictly inside a deleted interval. The
\emph{cut map} $\cmap$ sends a retained cut $p$ of $w$ to the number of
retained positions before it, a cut of $u$; it is monotone, and the two ends
of a deleted interval have the same image, a \emph{junction} of $u$. Every
cut $a$ of $u$ has a greatest preimage $\gamma(a)$, which is the right end
of a deleted interval if $a$ is a junction. If no junction lies strictly
between the cuts $a<a'$ of $u$, then
\begin{equation}\label{eq:between-junctions}
  \wsub ua{a'}=\wsub w{\gamma(a)}{\gamma(a)+a'-a},
\end{equation}
and $\gamma(a)+a'-a$ is the left end of the deleted interval of $a'$ if $a'$
is a junction and $\gamma(a')$ otherwise. If the two ends of every deleted
interval carry the same state, then the state of $u$ at $\cmap(p)$ is the
state of $w$ at $p$, for every retained cut $p$.

\subparagraph*{Geometry of the sites.}
\begin{enumerate}[label=\textup{(G\arabic*)}]
\item\label{it:G1} $s_o<c_o<c_o'<e_o$, since $1\le i<j<b$ and the parse of
  $o$ has at least $b$ lower blocks. Sites of different traversing
  occurrences lie in different upper blocks. So in $u_l$ the junction of $o$
  is followed by the $\lambda$ letters $z_\lambda$ up to $\cmap(e_o)$, the
  cuts strictly between the junction and $\cmap(e_o)$ are images of cuts
  strictly inside the upper block $o$, and two junctions of one word are
  more than $\lambda$ apart.
\item\label{it:G2} A word with a traversing occurrence has a level-$t$
  anchor before its end, so it is active at level $t+1$.
\item\label{it:G3} A level-$(t+1)$ marker ends at $e_o$, and every such
  marker has length at most $\lambda$: otherwise its occurrence would
  contain the $i+1$ lower anchors at the depths $i,\dots,0$, and
  $i+1>k(N-1)+2$, contradicting \cref{lem:depth-span}. For the same reason
  no site meets the last $\len{\mk{t+1}l}$ letters of a word $w_l$ with
  $l\in\Act_{t+1}$, so $u_l$ ends with $\mk{t+1}l$.
\end{enumerate}

\subparagraph*{States.}
Let $\omega$ be the common value $\omega(l,c_o)=\omega(l,c_o')$ of the equal
labels. If $\omega>0$, the same marker of length $\omega$ ends at both cuts,
so the $\omega$ letters before $c_o$ and before $c_o'$ are equal, and the
equal states at the distance $\omega$ give
\begin{equation}\label{eq:window}
  \wsub{w_l}{c_o-x}{c_o}=\wsub{w_l}{c_o'-x}{c_o'},\qquad
  q_l(c_o-x)=q_l(c_o'-x)\qquad\text{for }0\le x\le\omega.
\end{equation}
In particular the two ends of every site carry the same state. Hence the
states of $U$ at the images of retained cuts are the states of $W$: the
endpoints agree, no new state is visited, and $U$ is fresh, because a cut
before the end of $u_l$ is the image of a retained cut before the end of
$w_l$. The tuple $U$ is strictly shorter than $W$, because the occurrence
with the parse $Z$ traverses the edge.

\begin{lemma}[Anchors are preserved]\label{lem:anchors-preserved}
For every level $r$ with $1\le r\le t+1$ and every coordinate $l$, the
level-$r$ anchors of the text $u_l$, for the markers of $W$, are exactly the
images $\cmap(e)$ of the retained cuts $e$ of $w_l$ that are level-$r$
anchors of the text $w_l$.
\end{lemma}

\begin{proof}
We delete the sites of $w_l$ one at a time, from right to left, and show
that each deletion has the stated property. The following three facts about
a site $[c,c')$ concern only the prefix $\wpre{w_l}{c'}$, which the
deletions to its right do not change.
\begin{enumerate}[label=\textup{(V\arabic*)}]
\item\label{it:V1} If $t\ge1$, a level-$t$ marker ends at $c$, because $c$
  is a lower anchor other than the end of the word.
\item\label{it:V2} For $1\le r\le t$, a level-$r$ marker ends at $c$ if and
  only if it ends at $c'$, by the equal labels.
\item\label{it:V3} No level-$(t+1)$ marker ends at $c$ or at $c'$, by
  \ref{it:G1}.
\end{enumerate}
Write the current word as $x=yFv$ with $\len y=c$, and $x'=yv$. Consider an
occurrence of a level-$r$ marker $\mu$ in $x'$. If it lies in $y$ or in $v$
it is an occurrence in $x$ ending at a retained cut with the same image,
and conversely. If it ends at the junction, it is a suffix of $y$; by
\ref{it:V2} and \ref{it:V3} this is the case if and only if $\mu$ ends at
$c$ and at $c'$ in $x$. Now let $\mu=SR$ with $S$ a nonempty suffix of $y$
and $R$ a nonempty prefix of $v$. We claim $\len S\le\omega$. Otherwise
$t\ge1$, since for $t=0$ the markers are letters, and by \ref{it:V1} a
level-$t$ marker $\mu_t$ of length at most $\omega$ ends at $c$, so its
occurrence lies inside that of $\mu$ and ends before it. If $r\le t$, the
level-$r$ marker of the coordinate of $\mu_t$ is a suffix of $\mu_t$, so a
level-$r$ marker occurs inside the occurrence of $\mu$ and ends before it,
contradicting \cref{lem:markers}\eqref{it:mk-anti}. If $r=t+1$, then by
\cref{lem:markers}\eqref{it:mk-two} $\mu_t$ is the level-$t$ marker with
which $\mu$ begins, so $S=\mu_t$ has length at most $\omega$. This proves
the claim, and by \eqref{eq:window} $S$ is a suffix of $yF$, so $\mu$ ends
at the corresponding retained cut of $x$. Conversely, let $\mu$ end at a cut
after $c'$ in $x$ and begin before $c'$, and write $\mu=SR$ with $S$ a
suffix of $yF$ and $R$ a prefix of $v$. The same argument, with the
level-$t$ marker that ends at $c'$, gives $\len S\le\omega$, so $S$ is a
suffix of $y$ and $\mu$ ends at the image of the cut in $x'$. This covers
the occurrences that begin inside the site or before it.
\end{proof}

By \ref{it:G1} and \cref{lem:anchors-preserved}, no cut strictly between a
junction of $u_l$ and the cut $\lambda$ letters after it is a level-$(t+1)$
anchor of $u_l$.

\begin{lemma}[Transfer]\label{lem:transfer}
Let $a_2$ be a cut of $u_l$ such that $\wsub{u_l}{a_2}{a_2+\lambda}=
z_\lambda$ and no junction of $u_l$ lies strictly between $a_2$ and
$a_2+\lambda$. Then $a_2$ is a junction of $u_l$.
\end{lemma}

\begin{proof}
By \eqref{eq:between-junctions}, $\wsub{w_l}{c}{c+\lambda}=z_\lambda$ for
$c=\gamma(a_2)$. By \ref{it:G3} a level-$(t+1)$ marker of length at most
$\lambda$ is a suffix of $z_\lambda$, so $c+\lambda$ is an upper anchor of
$w_l$. By \cref{lem:traverse}, which applies since $i\ge a+\Marg$, the upper
occurrence ending at $c+\lambda$ traverses the edge. Its site is
$[c-\delta,c)$, so $a_2=\cmap(c)$ is a junction.
\end{proof}

\begin{proof}[Proof of \cref{thm:block-deletion}]
It remains to show $\Chan_{hg}(U)\subseteq\Chan_{hg}(W)$ for a host $h$ and
a target $g$. Let $a',b'$ be cuts of $u_h$ and $u_g$ with
$T=\wsuf{u_g}{b'}\pre\wsuf{u_h}{a'}$, and let $\theta=\len T$. We must find
cuts of $w_h$ and $w_g$ with the same states whose rests are compared. A
\emph{host junction offset} is a number $x$ with $0<x<\theta$ such that
$a'+x$ is a junction of $u_h$, and a \emph{target junction offset} one such
that $b'+x$ is a junction of $u_g$.

\emph{Step 1: junction offsets correspond.} Let $x$ be a target junction
offset. By \ref{it:G1}, $x+\lambda\le\theta$, the target reads $z_\lambda$
from $b'+x$, and so does the host from $a'+x$. By \ref{it:G3} a
level-$(t+1)$ marker is a suffix of $z_\lambda$, so $a'+x+\lambda$ is a
level-$(t+1)$ anchor of $u_h$. If a host junction offset $y$ satisfied
$x<y<x+\lambda$, this anchor would lie strictly between a junction and the
cut $\lambda$ letters after it, which is impossible. So \cref{lem:transfer}
applies, and $x$ is a host junction offset. Symmetrically, a host junction
offset $y$ with $y+\lambda\le\theta$ is a target junction offset: the target
reads $z_\lambda$ from $b'+y$, the cut $b'+y+\lambda$ is a level-$(t+1)$
anchor of $u_g$, and no target junction offset lies strictly between $y$ and
$y+\lambda$.

\emph{Step 2: the matched case.} Suppose every host junction offset is a
target junction offset, and let $x_1<\dots<x_m$ be the junction offsets. By
\eqref{eq:between-junctions}, applied between consecutive junctions,
$\wsuf{w_g}{\gamma(b')}$ is the word obtained from $T$ by inserting $F$ at
the offsets $x_1,\dots,x_m$, and the same word is a prefix of
$\wsuf{w_h}{\gamma(a')}$. The states at $\gamma(a')$ and $\gamma(b')$ are
those of $U$ at $a'$ and $b'$. This includes the case of an empty $T$.

\emph{Step 3: the unmatched case.} Otherwise, by Step 1, there is a host
junction offset $y$ with $y+\lambda>\theta$, and it is the only host
junction offset that is not a target junction offset, by \ref{it:G1}. We
first show that the target has no junction at or after $b'$. If
$g\notin\Act_{t+1}$, the word $w_g$ has no site, by \ref{it:G2}. Let
$g\in\Act_{t+1}$ and $\sigma=\mk{t+1}g$, a suffix of $u_g$ by \ref{it:G3}.
If $\len\sigma\le\theta$, then $\sigma$ ends at the cut $a'+\theta$ of
$u_h$, a level-$(t+1)$ anchor strictly between the junction $a'+y$ and the
cut $\lambda$ letters after it, which is impossible. So $\len\sigma>\theta$,
and by \ref{it:G3} no junction of $u_g$ lies at or after $b'$. Consequently
$T=\wsuf{w_g}{\tilde b}$ for $\tilde b=\len{w_g}-\theta$, with the state of
$b'$, there is no target junction offset, and by Step 1 $y$ is the only host
junction offset.

Let $o$ be the traversing occurrence of the junction $a'+y$, with the site
$[c,c')$, and write $T=SR$ with $\len S=y$; $R$ is not empty. By
\eqref{eq:between-junctions}, $S=\wsub{w_h}{c-y}{c}$, with the state of $a'$
at the cut $c-y$, and $R\pre\wsuf{w_h}{c'}$. If $y\le\omega$, then by
\eqref{eq:window} $S=\wsub{w_h}{c'-y}{c'}$ and the state at $c'-y$ is that
at $c-y$; so $\wsuf{w_g}{\tilde b}=SR\pre\wsuf{w_h}{c'-y}$ is a comparison
of $W$ with the required states.

It remains to exclude $y>\omega$. For $t=0$, the word $w_g$ is not empty, so
$g\in\Act_1$ and $\len{\mk1g}=1\le\theta$, which was excluded above. Let
$t\ge1$. By \ref{it:V1} a level-$t$ marker $\mu$ of length at most
$\omega<y$ ends at $c$; it is a suffix of $S$, so it occurs in $w_g$, begins
at or after $\tilde b$, and ends at the cut $\tilde b+y<\len{w_g}$. So $w_g$
has a level-$t$ anchor before its end, $g\in\Act_{t+1}$, and since
$\theta<\len\sigma$ this occurrence of $\mu$ lies inside the final
occurrence of $\sigma=\mk{t+1}g$, begins after its beginning and ends before
its end. This contradicts \cref{lem:markers}\eqref{it:mk-two}.
\end{proof}

\subsection{The Bound}\label{app:fresh-bound}

\begin{lemma}[Labels are local]\label{lem:labels-local}
Let $(Z,a,b)$ be an edge at level $t$ and $d$ a depth with $1\le d\le
b-\Marg$. There are marker texts $\tau_1,\dots,\tau_t$, some possibly
undefined, and a map $g_d\colon S\to S$ such that for every traversing upper
occurrence $o$ of a word $w_l$, the label of its cut at depth $d$ is
$(\tau_1,\dots,\tau_t,g_d(q_l(c)))$, where $c$ is the cut of $o$ at depth
$b$.
\end{lemma}

\begin{proof}
Let $y_d$ be the text of the lower blocks of $Z$ at the depths
$b-1,\dots,d$; it is the text between the cuts of $o$ at the depths $b$ and
$d$. A level-$r$ marker with $1\le r\le t$ that ends at the cut at depth $d$
contains at most $k(N-1)+2$ lower anchors, by \cref{lem:depth-span} for the
levels $r-1$ and $r$, because lower anchors are level-$(r-1)$ anchors. The
cuts at the depths $d,\dots,b$ are $b-d+1>k(N-1)+2$ lower anchors, so the
marker begins at or after $c$: it is a suffix of $y_d$. So $\tau_r$ is the
level-$r$ marker that is a suffix of $y_d$, and the state before the longest
of them is obtained from $q_l(c)$ by reading a prefix of $y_d$.
\end{proof}

\begin{proof}[Proof of \cref{thm:fresh-bound}]
Let $W$ be a minimal fresh tuple. By \cref{lem:depth-span} its top level
$s$ is at most $k(N-1)$. Fix $t\le s$. There are at most $(k+1)^t$
possibilities for the texts $\tau_r$ of \cref{lem:labels-local}, since at
each level at most one of at most $k$ marker texts ends at a cut, and
$N^N$ maps $g_d$. If an edge $(Z,a,b)$ had $b-a>E$, two of the
$b-a-2\Marg+1$ depths between $a+\Marg$ and $b-\Marg$ would have the same
texts and the same map, and by \cref{lem:labels-local} they would form a
block deletion, contradicting \cref{thm:block-deletion}. So every edge has
at most $E$ lower blocks.

Let $h_t$ be the number of distinct texts of upper blocks at level $t$ and
$P_t$ the number of distinct parses. By \cref{lem:parse-local}, the parse of
an upper block with the text $z$ is determined by $z$ and its set of
crossing anchors, a subset of the set $\operatorname{Pre}(z)$ of at most
$k(N-1)$ positions; so $P_t\le F\,h_t$. The path of a parse $Z$ has at most
$P_t+2$ nodes, the first is $0$, the last is the length of $Z$, and
consecutive nodes are at most $E$ apart; so $Z$ has at most $(P_t+2)E$ lower
blocks. For $t\ge1$, every upper block of level $t-1$ is a lower block in
the parse of an upper block of level $t$, hence
\[
  h_{t-1}\le P_t\,(P_t+2)\,E\le F h_t\,(F h_t+2)\,E .
\]
At the top level the upper blocks are the nonempty words, so $h_s\le k$. By
induction $h_{s-d}\le\nu_d$, and since the $\nu_d$ are nondecreasing and
$s\le kN$, every upper block of every level has at most $X=(\nu_{kN}F+2)E$
lower blocks. Blocks of level $0$ are letters, so a block of level $t+1$ has
at most $X^{t+1}$ letters, and each word, a block of level $s+1$ if it is not
empty, has at most $X^{kN+1}$ letters. The closed form is proved in
\cref{app:complexity-closed}.

The second statement of \cref{thm:fresh} follows because a fresh dominator
of least total length of a fresh tuple is a minimal fresh tuple.
\end{proof}
